A transformação homogênea entre os referenciais consecutivos $\{i-1\}$ e $\{i\}$ é obtida a partir da convenção de \gls{dh} \cite{craig2005}. 
Cada transformação descreve a relação geométrica entre dois referenciais adjacentes e depende dos quatro parâmetros apresentados na Tabela~\ref{tab:dh_zxx}.

A matriz de transformação ${}^{i-1}T_i$ é construída pela composição de quatro operações fundamentais: uma rotação de $\alpha_{i-1}$ em torno de $\hat{X}_{i-1}$, seguida de uma translação de $a_{i-1}$ ao longo de $\hat{X}_{i-1}$, uma rotação de $\theta_i$ em torno de $\hat{Z}_i$ e, por fim, uma translação de $d_i$ ao longo de $\hat{Z}_i$ \cite{craig2005}.
Assim, utilizando as matrizes elementares de rotação ($R$) e de translação ($D$), tem-se

\begin{equation}
{}^{i-1}T_i
=
R_X(\alpha_{i-1})\,D_X(a_{i-1})\,R_Z(\theta_i)\,D_Z(d_i),
\label{eq:Ti_1}
\end{equation}


onde $R_X(\alpha_{i-1})$ representa a rotação de ângulo $\alpha_{i-1}$ em torno de $\hat{X}_{i-1}$,
$D_X(a_{i-1})$ representa a translação de magnitude $a_{i-1}$ ao longo de $\hat{X}_{i-1}$,
$R_Z(\theta_i)$ representa a rotação de ângulo $\theta_i$ em torno de $\hat{Z}_i$ e
$D_Z(d_i)$ representa a translação de magnitude $d_i$ ao longo de $\hat{Z}_i$.


%%%%%%%%%%%%%%%%%%%%%%%%%%%%%%%%%%%

As quatro matrizes elementares utilizadas na convenção de \gls{dh} podem ser escritas como \cite{craig2005}:

\begin{equation}
R_X(\alpha_{i-1}) =
\begin{bmatrix}
1 & 0 & 0 & 0 \\
0 & \cos\alpha_{i-1} & -\sin\alpha_{i-1} & 0 \\
0 & \sin\alpha_{i-1} &  \cos\alpha_{i-1} & 0 \\
0 & 0 & 0 & 1
\end{bmatrix},
\end{equation}

\begin{equation}
D_X(a_{i-1}) =
\begin{bmatrix}
1 & 0 & 0 & a_{i-1} \\
0 & 1 & 0 & 0 \\
0 & 0 & 1 & 0 \\
0 & 0 & 0 & 1
\end{bmatrix},
\end{equation}

\begin{equation}
R_Z(\theta_i) =
\begin{bmatrix}
\cos\theta_i & -\sin\theta_i & 0 & 0 \\
\sin\theta_i &  \cos\theta_i & 0 & 0 \\
0 & 0 & 1 & 0 \\
0 & 0 & 0 & 1
\end{bmatrix},
\end{equation}

\begin{equation}
D_Z(d_i) =
\begin{bmatrix}
1 & 0 & 0 & 0 \\
0 & 1 & 0 & 0 \\
0 & 0 & 1 & d_i \\
0 & 0 & 0 & 1
\end{bmatrix}.
\end{equation}

Multiplicando essas quatro operações na ordem prescrita, conforme a equação \eqref{eq:Ti_1}, obtém-se a forma geral \cite{craig2005}:

\begin{equation}
\label{eq:transf_homogenea_geral_Craig}
{}^{i-1}T_i =
\begin{bmatrix}
\cos\theta_i & -\sin\theta_i & 0 & a_{i-1} \\
\sin\theta_i\cos\alpha_{i-1} &
\cos\theta_i\cos\alpha_{i-1} &
-\sin\alpha_{i-1} &
-\sin\alpha_{i-1}d_i \\
\sin\theta_i\sin\alpha_{i-1} &
\cos\theta_i\sin\alpha_{i-1} &
\cos\alpha_{i-1} &
\cos\alpha_{i-1}d_i \\
0 & 0 & 0 & 1
\end{bmatrix}.
\end{equation}


As transformações homogêneas representam, em uma única matriz $4\times 4$, tanto a orientação quanto a posição de um elo em relação ao anterior.
Cada matriz ${}^{i-1}T_i$ descreve a rotação e a translação do referencial $\{i\}$ em relação ao referencial $\{i-1\}$, constituindo a forma geral da transformação entre elos consecutivos de um manipulador.
Essa formulação é utilizada no cálculo direto da cinemática de posição e orientação do efetuador, permitindo expressar a pose final (posição e orientação) em função das variáveis articulares do manipulador.
